% To be reformatted in the TAC style (tac.cls) upon acceptance.
\documentclass[11pt]{amsart}

\usepackage{amsmath,amssymb}
\usepackage[hidelinks]{hyperref}
\usepackage{tikz}

\newtheorem{thm}{Theorem}[section]
\newtheorem{lem}[thm]{Lemma}
\newtheorem{cor}[thm]{Corollary}
\newtheorem{prop}[thm]{Proposition}
\theoremstyle{definition}
\newtheorem{defi}[thm]{Definition}
\newtheorem{rem}[thm]{Remark}
\newtheorem{obs}[thm]{Observation}

\newcommand{\I}{\mathbb{I}}
\newcommand{\DM}[1]{\mathrm{DM}(#1)}
\newcommand{\Dbox}{\square_{\mathrm{DM}}}
\newcommand{\aconv}{\equiv}
\newcommand{\nconv}{\not\equiv}
\newcommand{\refl}{\mathsf{refl}}
\newcommand{\Path}{\mathsf{Path}}
\newcommand{\PathP}{\mathsf{PathP}}
\newcommand{\U}{\mathcal{U}}
\newcommand{\hcomp}{\mathsf{hcomp}}
\newcommand{\act}[1]{A_{#1}}
\newcommand{\cellp}{\mathsf{p}}
\newcommand{\Real}[1]{|#1|}

\begin{document}

\title[Cellular contexts realize their cube presheaves]{Generic
  cellular contexts in cubical type theory\\ realize their cube
  presheaves on the nose}

\author{Ryota Shibusawa}
\address{Daiichi Institute of Technology, Japan}
\email{r-shibusawa@daiichi-koudai.ac.jp}

\subjclass[2020]{18N45, 03B38, 18M99}
\keywords{cubical type theory, cube categories, computads,
  polygraphs, connections, De Morgan algebra, definitional equality}

\begin{abstract}
  A companion paper showed that in cubical type theory with a De
  Morgan interval, a generic higher \emph{loop} spans, under
  definitional equality, a boundary-collapsed representable presheaf
  on the De Morgan cube category --- a ``De Morgan sphere'' --- and
  left open what happens for cells with \emph{generic} boundaries.
  This paper answers that question.  For a cellular context --- a
  cubical computad: generic vertices, generic edges between them,
  generic squares with prescribed edges, and so on --- the cubes
  definable by pure reparametrization form, modulo definitional
  equality and naturally in the cube dimension, exactly the
  \emph{cellular presheaf} of the computad: the colimit of one
  representable per cell, glued along the attachment maps.  The
  gluing is definitional in the strongest sense: each face of a
  generic cell \emph{is} its prescribed boundary cell on the nose,
  and no further identifications occur --- in particular, cells
  sharing a face are identified along exactly that face.  For a
  single cell with generic boundary the layer is the full
  representable, with no collapse; the spheres of the companion
  paper are recovered as the boundary-degenerate limit.  As a
  by-product we isolate a phenomenon absent from the Kan world:
  \emph{strict fillers are not unique} --- a boundary in the strict
  layer can admit several distinct fillers, reflecting the
  non-injectivity of the restriction of free De Morgan polynomials
  to the faces of a cube.  All positive and negative instances are
  machine-checked in a self-contained proof-assistant kernel, and
  every statement is a decidable, assumption-free fact about the
  algorithmic equality of a minimal fragment.
\end{abstract}

\maketitle

\section{Introduction}
\label{sec:intro}

Cubical type theories \cite{CCHM} equip the identity structure of
types with definitional operations indexed by a De Morgan algebra:
reparametrizing a cell along $\wedge$, $\vee$, $\neg$ --- the
connections and reversals --- happens on the nose, not up to higher
cells.  A companion paper \cite{Sphere} determined the exact extent
of this strictness over contexts of generic \emph{loops}: the
definable cubes form a wedge of \emph{De Morgan spheres},
boundary-collapsed representable presheaves $\Dbox(-,[n])/\partial$
on the De Morgan cube category $\Dbox$, one per loop, and nothing
else is identified.  The collapse of the boundary was forced by the
choice of generic objects: a loop's faces are reflexivities, and
reflexivities absorb everything applied to them.

That paper left open the \emph{unpointed} question: what does a
generic cell with a \emph{generic} boundary span?  The present
paper answers it, and the answer upgrades the picture from wedges of
quotients to genuine cell complexes.

\subparagraph*{Cellular contexts.}  By a \emph{cellular context} (a
\emph{cubical computad}) we mean a context of generic cells
declared dimension by dimension: generic points $a, b, \dots : A$;
generic edges between prescribed points,
$p : \Path\,A\,a\,b$; generic squares with four prescribed edges,
$Q : \PathP\,(\langle i \rangle\, \Path\,A\,(r\,i)\,(s\,i))\,p\,q$;
and so on.  This is the type-theoretic incarnation of a computad
(polygraph) \cite{Burroni,Polygraphs} over the cube category: free
generators in each dimension, attached along specified lower cells.
Every computad $C$ has a \emph{cellular presheaf} $\Real{C}$ on
$\Dbox$ --- one representable $\Dbox(-,[n_c])$ per cell $c$, glued
along the attachment maps --- the cubical analogue of the cell
complex a CW-presentation describes.

\subparagraph*{Main results.}  Let $\aconv$ denote the definitional
(algorithmic) equality of the calculus
(Section~\ref{sec:calc}).

\begin{enumerate}
\item \emph{Attachment is definitional}
  (Section~\ref{sec:attach}).  Each face of a generic cell computes,
  on the nose, to its prescribed boundary cell:
  $\langle j \rangle\, Q(0)(j) \aconv p$, \
  $\langle i \rangle\, Q(i)(0) \aconv r$, and so on; corners
  compute through two steps to the prescribed vertices.  The
  structure maps of the cellular presheaf are thus judgmental
  equalities, not merely admissible ones.
\item \emph{The realization theorem}
  (Sections~\ref{sec:single}--\ref{sec:main}).  The
  pure-reparametrization cubes over a cellular context, modulo
  $\aconv$ and naturally in the cube dimension $[m]$, are exactly
  the cellular presheaf:
  \[
    \{\text{$m$-cubes}\}/{\aconv} \;\cong\; \Real{C}([m]).
  \]
  For a single cell with generic boundary this is the \emph{full}
  representable $\Dbox(-,[n])$ --- no collapse; for cells sharing a
  face, the gluing identifies exactly the shared face
  (machine-checked in both directions: the cross-cell face equation
  holds definitionally, and interior cubes of distinct cells with
  \emph{identical boundaries} remain separated).
\item \emph{Strict fillers are not unique}
  (Section~\ref{sec:fillers}).  The elements
  $i \wedge \neg i$ and $(i \wedge \neg i) \wedge (j \vee \neg j)$
  of the free De Morgan algebra are distinct, yet restrict equally
  to all four faces of the square.  Consequently two
  \emph{different} strict squares can fill the \emph{same} boundary.
  Where Kan filling is unique up to homotopy, strict filling is
  non-unique on the nose --- for the entirely algebraic reason that
  restriction of free De Morgan polynomials to the boundary of a
  cube is not injective.
\item \emph{Sphere recovery} (Section~\ref{sec:spheres}).
  Degenerating the boundary of the computad to a point recovers the
  wedge-of-spheres theorem of \cite{Sphere}: the collapse
  $\Dbox(-,[n]) \twoheadrightarrow \Dbox(-,[n])/\partial$ is the
  image, under realization, of collapsing the boundary cells.
\end{enumerate}

All positive and negative instances are elaboration-time checks in
the artifact kernel (Section~\ref{sec:artifact}); by the
decidability theorem imported from the first companion paper
\cite{Const}, each is a decided, assumption-free fact.

\subparagraph*{Honest positioning.}  As in \cite{Sphere}, the
positive direction --- that the De Morgan laws and the attachment
equations hold --- is by design of the theory.  The content is the
exactness: the identifications are \emph{exactly} those of the
cellular presheaf, no more (the separations, including the delicate
same-boundary ones behind non-uniqueness of fillers) and no less
(the attachment equations).  We know of no prior statement of the
realization theorem or of the non-uniqueness phenomenon.  The
proofs are elementary given two small kernel facts (K1, K2 below)
and are designed to be checked against the artifact.

\subparagraph*{Relation to the companions.}  This paper is the
third in a series but is written to be read independently.  From
\cite{Const} we import only the decidability of the algorithmic
equality on the minimal fragment (K3) and the two kernel facts K1,
K2; from \cite{Sphere} we import the De Morgan cube category and
the sphere theorem, which the present results generalize and
recover.  Nothing here depends on the transport-layer analysis of
\cite{Const} (Remark~\ref{rem:robust}).

\section{The De Morgan cube category, computads, and cellular
  presheaves}
\label{sec:cube}

\begin{defi}[{\cite[Def.~2.2]{Sphere}}]\label{def:cube}
  $\Dbox$ has objects $[m]$, $m \ge 0$, and morphisms
  $\Dbox([m],[n]) := \DM{i_1,\dots,i_m}^{\,n}$, tuples of free De
  Morgan polynomials, composed by substitution; the identity of
  $[n]$ is the generator tuple.  $\DM{\vec i}$ is the free De
  Morgan algebra (no excluded middle), with canonical antichain
  disjunctive normal forms; $0, 1$ are its \emph{endpoints}.  The
  \emph{$k$-th face} is $\delta_k^\varepsilon := (i_1, \dots,
  \varepsilon, \dots, i_{n-1}) : [n-1] \to [n]$; a morphism is
  \emph{face-factoring} if it factors through some
  $\delta_k^\varepsilon$, equivalently iff some coordinate is an
  endpoint \emph{of the free algebra}.
\end{defi}

\begin{defi}[Cellular contexts]\label{def:computad}
  A \emph{cellular context of dimension $\le 2$} over a base type
  $A : \U$ is a context declaring, in order:
  \begin{itemize}
  \item generic vertices $a_1, \dots : A$;
  \item generic edges $e : \Path\,A\,a\,a'$ between prescribed
    vertices;
  \item generic squares
    $Q : \PathP\,(\langle i \rangle\,
    \Path\,A\,(r\,i)\,(s\,i))\,p\,q$ with four prescribed edges
    $p, q, r, s$ whose endpoints match at the corners.
  \end{itemize}
  Higher cells are declared analogously along the cube-boundary
  telescopes; we spell out dimension $\le 2$, which is where our
  machine evidence lives, and treat higher dimensions uniformly
  (Remark~\ref{rem:highdim}).  We call the data a \emph{cubical
  computad} $C$; it is precisely a polygraph
  \cite{Burroni,Polygraphs} over $\Dbox$ presented type-theoretically.
\end{defi}

\begin{defi}[Cellular presheaf]\label{def:cellpresheaf}
  The \emph{cellular presheaf} $\Real{C}$ of a computad $C$ is the
  presheaf on $\Dbox$ with
  \[
    \Real{C}([m]) \;:=\;
    \Bigl(\coprod_{c \in C} \Dbox([m],[n_c])\Bigr)\Big/\sim,
  \]
  where $\sim$ is generated by the attachments: if
  $\vec\varphi = \delta_k^\varepsilon \circ \vec\chi$ face-factors,
  then $(c, \vec\varphi) \sim (\partial_k^\varepsilon c, \vec\chi)$,
  for $\partial_k^\varepsilon c$ the prescribed boundary cell.
  Equivalently, $\Real{C}$ is the colimit of the representables
  $\Dbox(-,[n_c])$ along the face inclusions specified by the
  attachments --- the cubical cell complex presented by $C$.
\end{defi}

\section{The calculus and imported facts}
\label{sec:calc}

We work in the reparametrization fragment of a CCHM-style cubical
type theory \cite{CCHM}: path types over a base type, path
abstraction and application; Kan operations are excluded from the
layer under study.  Definitional equality is the algorithmic
equality $\aconv$ of the concrete normalization-by-evaluation
kernel of \cite{Const}, whose two load-bearing clauses are:

\begin{enumerate}
\item[(K1)] \emph{Endpoint collapse is normal-form exact}: applying
  a path value at $r \in \DM{\vec i}$ collapses to the stored
  endpoint iff $r$ is an endpoint of the free algebra (the kernel
  decides by antichain normal forms).  For a value of a
  \emph{dependent} path type $\PathP\,(\langle i \rangle
  T)\,l\,u$, the stored endpoints are $l, u$; for a neutral spine,
  the stored endpoint \emph{annotations} of its type.
\item[(K2)] \emph{Conversion of neutrals is spine-wise}: same head
  variable, same length, convertible endpoint annotations, interval
  arguments equal in the free algebra.
\item[(K3)] \emph{Decidability, assumption-free} \cite{Const}: on
  the minimal fragment, evaluation and conversion terminate and
  $\aconv$ is decidable.
\end{enumerate}

\begin{rem}[Robustness]\label{rem:robust}
  The fragment involves no transport, so all results are
  insensitive to the transport-layer design of \cite{Const} and
  hold verbatim for a standard CCHM-style conversion; we use the
  kernel of \cite{Const} because there they are unconditional
  decidable facts (K3).
\end{rem}

Throughout, fix a cellular context; for a generic square $Q$ with
edges $p, q, r, s$ (as in Definition~\ref{def:computad}) and
$\varphi, \psi \in \DM{i,j}$ write
\[
  Q(\varphi)(\psi)
\]
for the doubly applied cell (with the boundary-typed endpoint
annotations) and
$\act{Q}(\vec\varphi) := \langle \vec i\, \rangle\,
Q(\varphi_1)(\varphi_2)$ for the reparametrized cube, and likewise
for edges.

\section{Attachment is definitional}
\label{sec:attach}

\begin{thm}[Attachment equations]\label{thm:attach}
  For a generic square $Q$ with prescribed edges $p, q, r, s$:
  \[
    \langle j \rangle\, Q(0)(j) \aconv p, \quad
    \langle j \rangle\, Q(1)(j) \aconv q, \quad
    \langle i \rangle\, Q(i)(0) \aconv r, \quad
    \langle i \rangle\, Q(i)(1) \aconv s,
  \]
  and corners compute through two steps:
  $\langle i \rangle \langle j \rangle\, Q(0)(0) \aconv
  \langle i \rangle \langle j \rangle\, a_{00}$, etc.
  More generally the whole face square is the edge square:
  $\langle i \rangle \langle j \rangle\, Q(0)(j) \aconv
  \langle i \rangle \langle j \rangle\, p(j)$.
\end{thm}
\begin{proof}
  By K1 at the outer application: $Q(0)$ collapses to the stored
  endpoint of the $\PathP$ type, which is the \emph{prescribed edge}
  $p$ --- a generic cell, not a reflexivity; the remaining
  application and $\eta$ give the first equation.  For the third,
  $Q(\varphi)$ at non-endpoint $\varphi$ is a neutral spine whose
  stored endpoint annotations are $r(\varphi), s(\varphi)$; K1 at
  the inner application returns them, and $\eta$ finishes.  Corners
  iterate the two collapses.  Machine probes: \texttt{gbAttPD},
  \texttt{gbAttQD}, \texttt{gbAttRD}, \texttt{gbAttSD},
  \texttt{gbCornerD}, \texttt{gbEdgeSqD}.
\end{proof}

\begin{rem}
  This is the exact point at which the generic-boundary situation
  departs from the loops of \cite{Sphere}: there the stored
  endpoints were reflexivity towers, through which every further
  application computed to the base point --- producing the
  $/\partial$ collapse.  Here the stored endpoints are generic
  cells, and the face-factoring morphisms land in the boundary's
  own layer instead of a single point.  The De Morgan laws
  themselves remain strict in the presence of a generic boundary:
  the unit and absorption laws hold at the full dependent square
  type (machine probes \texttt{gbEtaD}, \texttt{gbAbsorbD},
  \texttt{gbAbsorbPerturbD}).
\end{rem}

\section{The single-cell realization}
\label{sec:single}

\begin{lem}[Value computation, generic boundary]\label{lem:value}
  Over the square computad (four vertices, four edges, one square),
  open the $m$ binders of a definable cube at generic dimensions.
  The body evaluates to exactly one of:
  \begin{enumerate}
  \item a vertex variable;
  \item an edge spine $e(\chi)$ with $e \in \{p,q,r,s\}$ and $\chi$
    a non-endpoint;
  \item a full spine $Q(\varphi)(\psi)$ with both coordinates
    non-endpoints.
  \end{enumerate}
  A tuple $(\varphi,\psi)$ face-factoring as
  $\delta^\varepsilon \circ \chi$ evaluates to the \emph{same
  value} as the corresponding boundary composite: the edge spine of
  the prescribed face at $\chi$ (or a vertex, when $\chi$ is again
  an endpoint).
\end{lem}
\begin{proof}
  K1 case analysis as in Theorem~\ref{thm:attach}; the three cases
  are exhaustive because collapse consumes coordinates until either
  none remain (vertex) or the first non-endpoint argument freezes
  the spine.
\end{proof}

\begin{thm}[Single cell, full representable]\label{thm:single}
  The pure-reparametrization $m$-cubes over the square computad,
  modulo $\aconv$, are naturally isomorphic to
  $\Real{C} = \Dbox(-,[2])$ --- the \emph{full} representable, the
  isomorphism sending $\vec\varphi$ to $\act{Q}(\vec\varphi)$ and
  the face-factoring morphisms to the boundary composites.
\end{thm}
\begin{proof}
  \emph{Surjectivity} is Lemma~\ref{lem:value}: every value is a
  vertex, edge, or full spine, i.e.\ the image of a morphism (for
  the square computad, every boundary composite is a face-factoring
  morphism of $[2]$, so $\Real{C}([m]) = \Dbox([m],[2])$ ---
  attachments here \emph{create} no quotient, they only relate the
  coproduct's summands, and every summand embeds by face
  composition).

  \emph{Well-definedness} (the presheaf identifications hold) is
  Theorem~\ref{thm:attach}.

  \emph{Injectivity}: two full spines are convertible iff their
  tuples agree (K2, interval arguments componentwise in the free
  algebra); two edge spines iff same edge and same argument (K2:
  the four edges are distinct head variables); vertex values are
  the four distinct vertex variables; and the three classes are
  mutually non-convertible (distinct heads; a variable is not a
  spine).  A face-factoring tuple and a non-face-factoring tuple
  land in different classes by Lemma~\ref{lem:value}.  Finally two
  face-factoring tuples with the same value have the same boundary
  composite, hence are equal in $\Dbox([m],[2])$ --- e.g.\
  $p(\chi) = p(\chi')$ forces $\chi = \chi'$, and cross-edge
  coincidences occur only at corners, where the tuples agree.
\end{proof}

\begin{rem}[Higher dimensions]\label{rem:highdim}
  For a generic $n$-cell attached along the full cube-boundary
  computad of dimension $n$, Lemma~\ref{lem:value} and its
  consequences go through verbatim by induction on $n$: K1 peels
  endpoint coordinates one at a time, each collapse landing in the
  boundary computad one dimension down, and K2 separates the
  resulting spines by head and length.  We have carried out the
  machine verification for $n \le 2$; the induction step involves
  no new phenomenon, and we state the general case as a theorem
  with this uniform proof.
\end{rem}

\section{Gluing, and the main theorem}
\label{sec:main}

The induction step for several cells is the two-square computad:
squares $Q_1$ (edges $p, m, r_1, s_1$) and $Q_2$ (edges
$m, q, r_2, s_2$) \emph{sharing} the edge $m$.

\begin{thm}[Gluing along a shared face]\label{thm:glue}
  \begin{enumerate}
  \item (Identification.)  The two faces agree across the two
    cells, definitionally:
    \[
      \langle i \rangle \langle j \rangle\, Q_1(1)(j)
      \;\aconv\;
      \langle i \rangle \langle j \rangle\, Q_2(0)(j)
      \qquad (\text{both} \aconv \text{the edge square of } m).
    \]
  \item (No over-identification.)  Let
    $E_0 := i \vee \neg i \vee j \vee \neg j$ and
    $D_0 := i \wedge \neg i \wedge j \wedge \neg j$ --- the
    edge-one and edge-zero elements of $\DM{i,j}$.  The squares
    \[
      S_1 := \langle i \rangle \langle j \rangle\, Q_1(E_0)(j),
      \qquad
      S_2 := \langle i \rangle \langle j \rangle\, Q_2(D_0)(j)
    \]
    have their \emph{entire boundaries} on $m$: together with
    the edge square $\langle i \rangle \langle j \rangle\, m(j)$,
    both inhabit the type
    $\Path\,(\Path\,A\,a_{10}\,a_{11})\,m\,m$.  All three are
    pairwise non-convertible.
  \end{enumerate}
\end{thm}
\begin{proof}
  (1) Both sides collapse by K1 to the shared cell $m$
  (Theorem~\ref{thm:attach} applied in each square), hence to each
  other.  Machine probes \texttt{seAttM1D}, \texttt{seAttM2D},
  \texttt{seGlueD}.  (2) $E_0$ and $D_0$ are non-endpoints of the
  free algebra whose restrictions to all four faces are $1$
  resp.\ $0$; so $S_1, S_2$ are full spines with heads $Q_1$
  resp.\ $Q_2$, the edge square is an $m$-spine, and K2 separates
  the three by head.  Well-typedness at the common type is the
  boundary computation of (1).  Machine probes
  \texttt{seS1CtrlD}, \texttt{seS2CtrlD}, \texttt{seSmCtrlD} and
  the three separation guards.
\end{proof}

\begin{thm}[Realization]\label{thm:main}
  For every cubical computad $C$ (of dimension $\le 2$; higher
  dimensions per Remark~\ref{rem:highdim}), the
  pure-reparametrization $m$-cubes over the context of $C$, modulo
  the algorithmic equality, form naturally in $[m]$ the cellular
  presheaf:
  \[
    \{\text{$m$-cubes}\}/{\aconv} \;\cong\; \Real{C}([m]).
  \]
  All of this is decidable and assumption-free by K3.
\end{thm}
\begin{proof}
  Induction on the cells of $C$.  Adding a vertex or an edge is the
  base analysis of Lemma~\ref{lem:value} (an edge over its two
  vertices is the one-dimensional instance of
  Theorem~\ref{thm:single}).  Adding a square $Q$ attached along
  prescribed edges: the new definable bodies are exactly the
  $Q$-spines (Lemma~\ref{lem:value}, whose proof is local to the
  cell and its boundary); the attachment identifications hold
  definitionally (Theorem~\ref{thm:attach}), matching the
  presheaf's $\sim$; and no new identifications occur --- between a
  $Q$-spine and anything previously definable (K2, head
  separation; the same-boundary cases are exactly those of
  Theorem~\ref{thm:glue}(2), which is the general no-merge
  computation: the boundary-collapsed interior of a new cell never
  merges with old material or with another cell's interior).
  Naturality in $[m]$ --- compatibility with reparametrization ---
  is the strict functoriality of the actions, proved in
  \cite[Lemmas 4.2--4.3]{Sphere} and re-checked at dependent
  square types here (\texttt{gbEtaD}, \texttt{gbAbsorbD}).
\end{proof}

\section{Strict fillers are not unique}
\label{sec:fillers}

In the Kan world, a boundary determines its fillers up to homotopy.
In the strict layer the situation is sharper in one way --- fillers
are literal --- and looser in another:

\begin{thm}[Non-uniqueness of strict fillers]\label{thm:fillers}
  Over the square computad there exist two non-convertible squares
  with the \emph{same} boundary: with
  $\varphi_1 := i \wedge \neg i$ and
  $\varphi_2 := (i \wedge \neg i) \wedge (j \vee \neg j)$,
  \[
    \langle i \rangle \langle j \rangle\, Q(\varphi_1)(j)
    \;\nconv\;
    \langle i \rangle \langle j \rangle\, Q(\varphi_2)(j),
  \]
  yet both inhabit
  $\PathP\,(\langle i \rangle\,
  \Path\,A\,(r(i \wedge \neg i))\,(s(i \wedge \neg i)))\,p\,p$.
\end{thm}
\begin{proof}
  $\varphi_1 \ne \varphi_2$ in the free algebra: absorption fails
  because the clause $\{i, \neg i\}$ contains neither $\{j\}$ nor
  $\{\neg j\}$, so
  $(i \wedge \neg i) \wedge (j \vee \neg j) \ne i \wedge \neg i$.
  Yet their restrictions to the four faces coincide ($0$ on the
  $i$-faces, $i \wedge \neg i$ on the $j$-faces), so the two
  squares have convertible boundaries; being non-degenerate full
  spines with distinct interval arguments, K2 separates them.
  Machine probes \texttt{gbFill1D}, \texttt{gbFill2D} and the
  separation guard.
\end{proof}

\begin{prop}[The minimal perturbation is absorbed]
  \label{prop:absorbed}
  The edge-vanishing elements of $\DM{i,j}$ are exactly $0$ and
  $D_0 = i \wedge \neg i \wedge j \wedge \neg j$, and $D_0 \le
  \varphi$ for \emph{every} $\varphi \ne 0$ (the clause of $D_0$
  contains every literal, hence contains every clause).
  Consequently join-perturbation by $D_0$ never produces a new
  filler: $\varphi \vee D_0 = \varphi$
  (machine probe \texttt{gbAbsorbPerturbD}).  The non-uniqueness of
  Theorem~\ref{thm:fillers} therefore arises only through
  meet-constructions such as $\varphi_2$, i.e.\ through the
  non-injectivity of the boundary-restriction map
  $\DM{i,j} \to \DM{j}^2 \times \DM{i}^2$ --- an intrinsic feature
  of free De Morgan algebras, faithfully reflected by definitional
  equality.
\end{prop}

\begin{rem}
  Theorem~\ref{thm:fillers} is consistent with the realization
  theorem --- the representable $\Dbox(-,[2])$ genuinely has
  several elements over a single boundary configuration --- but it
  is worth isolating because it contradicts a natural first guess
  (that strict fillers, when they exist, are unique) and because it
  gives the boundary-restriction non-injectivity of free De Morgan
  algebras a geometric meaning inside type theory.
\end{rem}

\section{Recovering the spheres}
\label{sec:spheres}

The companion's generic loops are the boundary-degenerate limit of
the present setting: substituting the constant computad (one
vertex, reflexivity boundaries) for the boundary cells sends the
square computad to the $2$-loop context of \cite{Sphere}.  Under
this substitution the stored endpoints of the spine collapse to
reflexivity towers, through which K1 computes to the base point ---
precisely the mechanism that turns the full representable into the
boundary-collapsed one:

\begin{prop}\label{prop:spheres}
  Boundary degeneration realizes the quotient: the substitution of
  the point computad for the boundary carries the isomorphism of
  Theorem~\ref{thm:single} to the sphere isomorphism
  $\{\text{cubes}\}/{\aconv} \cong \Dbox(-,[2])/\partial$ of
  \cite{Sphere}, the induced map on presheaves being the canonical
  collapse $\Dbox(-,[2]) \twoheadrightarrow
  \Dbox(-,[2])/\partial$.
\end{prop}
\begin{proof}
  The substitution sends attachment values (edge spines, vertices)
  to reflexivity towers and the base point; Lemma~\ref{lem:value}
  degenerates clause-wise to the value computation of
  \cite[Lemma 5.1]{Sphere}, identifying exactly the face-factoring
  classes.
\end{proof}

Together, Theorems~\ref{thm:single}, \ref{thm:main} and
Proposition~\ref{prop:spheres} give a uniform picture: \emph{the
definitional equality of cubical type theory realizes cubical
computads as their cell complexes, on the nose; collapsing
boundaries collapses realizations accordingly}.

\section{Machine verification}
\label{sec:artifact}

All instances quoted above are elaboration-time checks in the
self-contained Lean~4 cubical kernel of \cite{Const}: repository
\begin{center}
  \url{https://github.com/r-shibusawa/cubical-strict-j}.
\end{center}
The probes
of this paper are the files \texttt{LibGenBoundary.lean}
(attachments, corners, edge squares, unit and absorption at
dependent square types, filler non-uniqueness) and
\texttt{LibSharedEdge.lean} (cross-cell gluing and the no-merge
separations), with the supporting proof notes
\texttt{docs/\allowbreak GenericBoundary.md}.  They are included
from release \texttt{v1.3.0}, archived at DOI
\href{https://doi.org/10.5281/zenodo.21638976}%
{\texttt{10.5281/zenodo.\allowbreak TODOD}}.  A successful \texttt{lake build}
replays every check; by K3 each check is a decision.  As in the
companions, the kernel supplies machine-checked \emph{instances};
the theorems are pen-and-paper results about the conversion
algorithm, whose two load-bearing clauses (K1, K2) are small
inspectable functions.

\section{Related work and outlook}
\label{sec:related}

Computads/polygraphs as presentations of free higher structures go
back to Street and Burroni \cite{Burroni}; see \cite{Polygraphs}
for a modern account.  The present results interpret the
\emph{free} structure they present inside cubical type theory
\cite{CCHM} literally: definitional equality itself realizes the
presented cell complex over the De Morgan cube category
\cite{GrandisMauri,BuchholtzMorehouse}, with nothing added and
nothing collapsed.  The companion papers \cite{Const,Sphere}
supply the metatheoretic platform (decidability of the algorithmic
equality on the fragment) and the pointed special case (loops and
spheres).

Two directions seem worthwhile.  First, \emph{mixed layers}: our
fragment excludes the Kan operations; over a cellular context, the
full theory's cubes include composites and fillers, and the
inclusion of the strict layer into the full layer is a
type-theoretic incarnation of the unit ``free strict
$\to$ free weak'' comparison, whose exact structure (in particular
which weak cells are detected by strict boundaries, given
Theorem~\ref{thm:fillers}) deserves study.  Second,
\emph{presentations}: since realization is exact, one can read off
normal forms for the free cubical $\omega$-category on a computad
in the De Morgan signature from the type theory's normal forms ---
suggesting a proof-assistant-based approach to word problems for
cubical polygraphs.

\subparagraph*{Acknowledgements.}  The metatheoretic platform and
the companion results were developed in \cite{Const,Sphere}.

\begin{thebibliography}{10}

\bibitem{CCHM}
C.~Cohen, T.~Coquand, S.~Huber, A.~M\"ortberg.
\newblock Cubical type theory: a constructive interpretation of the
  univalence axiom.
\newblock \emph{TYPES 2015}, LIPIcs 69, 2018.

\bibitem{GrandisMauri}
M.~Grandis, L.~Mauri.
\newblock Cubical sets and their site.
\newblock \emph{Theory and Applications of Categories} 11 (2003),
  185--211.

\bibitem{BuchholtzMorehouse}
U.~Buchholtz, E.~Morehouse.
\newblock Varieties of cubical sets.
\newblock \emph{RAMiCS 2017}, LNCS 10226; arXiv:1701.08189.

\bibitem{Burroni}
A.~Burroni.
\newblock Higher-dimensional word problems with applications to
  equational logic.
\newblock \emph{Theoretical Computer Science} 115 (1993), 43--62.

\bibitem{Polygraphs}
D.~Ara, A.~Burroni, Y.~Guiraud, P.~Malbos, F.~M\'etayer,
S.~Mimram.
\newblock \emph{Polygraphs: From Rewriting to Higher Categories}.
\newblock London Mathematical Society Lecture Note Series,
  Cambridge University Press, 2025; arXiv:2312.00429.

\bibitem{Const}
R.~Shibusawa.
\newblock Evaluation-time constancy inference for transport in
  cubical type theory.
\newblock Preprint, 2026.  Artifact:
  \url{https://github.com/r-shibusawa/cubical-strict-j}, archived
  at DOI \href{https://doi.org/10.5281/zenodo.21637898}%
  {10.5281/zenodo.21637898}.

\bibitem{Sphere}
R.~Shibusawa.
\newblock The strict layer of generic cells in cubical type theory
  is a wedge of De Morgan spheres.
\newblock Preprint, 2026.  Included in the artifact repository
  above (release \texttt{v1.3.0}).

\end{thebibliography}

\end{document}
